\documentclass[11pt]{article}
\textwidth 15.9cm
\textheight 23. cm
\addtolength{\oddsidemargin}{-11.5mm}
\addtolength{\evensidemargin}{-11.5mm}
\addtolength{\topmargin}{-19.0mm}

\usepackage[ansinew]{inputenc}
\usepackage{amsmath,amssymb}
\usepackage{bbm}
\usepackage{graphicx}
\usepackage{hyperref}

\include{GempicLatexMacros}

\renewcommand{\part}{\texttt{part}}
\newcommand{\cell}{\texttt{cell}}
\newcommand{\node}{\texttt{node}}
\newcommand{\prim}{\texttt{prim}}

\begin{document}


Currently only a bilinear filter is available (See \ref{sec:FilterInput} for the controls).

\section{Bilinear Filter}\label{sec:BilinearFilter}
The multi-pass bilinear filter (cf.~\cite{vay2011}) interpolates the input using $N$ one-dimensional stencils, where $N$ is the dimension of the simulation. The stencils themselves can be iteratively calculated by \lstinline{nPass} convolutions with themselves according to the formula:
\begin{equation}
    f_{new}[i] = \alpha f_{old}[i] + (1 - \alpha)\frac{f_{old}[i-1] + f_{old}[i+1]}{2},
\end{equation}
where $f_{old}[i]$ is the $i$'th index of the old filter, $f_{new}[i]$ the $i$'th index of the new filter, and $\alpha$ is the weighting parameters. For the bilinear filter, $\alpha=0.5$.
If \lstinline{compensation} was set in the input file, we perform a final iteration with $\alpha=\frac{n_{Pass}}{2}$.
In 2D, with two stencils $f_x$, $f_y$ the application of the filter is then
\begin{multline}
    out[i,j] = \sum_{i_f=0}^{f_{x_{len}}}\sum_{j_f=0}^{f_{y_{len}}} f_x[i_f]f_y[j_f]\big(in[i-i_f,j-j_f] + in[i-i_f,j+j_f]\\ + in[i+i_f,j-j_f] + in[i+i_f,j+j_f]\big)
\end{multline}
\end{document}